\documentclass[11pt]{scrartcl}
\usepackage[sexy]{evan} %BLBLBLBLBLBLBLBLBLBLBLB EVANCHEN
\usepackage[T1]{fontenc}
\usepackage[utf8]{inputenc}
\usepackage{graphicx}
\usepackage{amsmath}
\usepackage{tikz}
\usepackage{asymptote}
\usepackage{amsthm}
\usepackage{gensymb}
\usepackage{amsfonts}
\usepackage[english,italian]{babel}
\usepackage{quoting}
\quotingsetup{font=big}
\usepackage{booktabs}
\usepackage{caption}
\usepackage{lmodern}
\newcommand{\dif}{\text{d}}
\renewcommand{\kg}{\,\text{Kg}}

\begin{document}
	\title{Appunti di Fisica 1}
	\subtitle{UNIPD - Gianguido Dall'Agata}
	\date{Febbraio - Giugno 2026}
	
	
	\maketitle
	
	\tableofcontents
	
	\newpage
	
	\section{Concetti introduttivi}
	\subsection{Analisi dimensionale}
	Ogni grandezza in fisica ha una sua unità di misura. In particolare le tre unità di misura fondamentali sono 
	\begin{itemize}
		\item Massa (chilogrammo).
		\item Lunghezza (metro).
		\item Tempo (secondo).
	\end{itemize}
	Le altre grandezze sono tutte espresse come prodotti di potenze (possibilmente negative) di queste tre. La regola fondamentale in fisica è che \textbf{le unità di misura non si sommano}, quando sommiamo due grandezze queste devono sempre avere la stessa unità di misura.
	\subsection{Spazio, Tempo e sistemi di riferimento}
	Dato che questo sarà un concetto di meccanica newtoniana, ci affidiamo al concetto newtoniano di spazio e tempo. Per Newton questi sono concetti assoluti e primitivi, esiste un tempo e uno spazio completamente indipendente da chi lo osserva, e questo è intuitivo e non ha bisogno di essere chiarito. \\
	In realtà questa cosa è in diretta contraddizione con il principio di relatività galileiano, ma è una questione abbastanza vessata, e non entreremo in dettaglio. \\
	Per descrivere tempo e spazio abbiamo bisogno di un sistema di riferimento. In questo senso ci basiamo sulla trattazione del corso di geometria $1$, sappiamo cos'è un sistema di riferimento, e ne faremo ampio uso durante tutto il corso. Il sistema di riferimento standard è quello centrato nello $\vec{0}$ con $3$ assi "destrorsi" ortogonali tra loro e unitari.
	\subsection{Vettori}
	Perché lavoriamo con i vettori? Perché ci permettono di esprimere in modo conciso tante informazioni su poche grandezze. Le grandezze "vettoriali" sono quelle che hanno bisogno di essere descritte tramite un verso, una direzione e un modulo, e quindi per questo ci viene in aiuto la rappresentazione geometrica di vettore. Una grandezza espressa solo da un numero invece si dice grandezza scalare.
\end{document}
